\chapter{Barrick Mining Corporation and Valuation Perimeter}
\chapterstatus{Legal identity and listing continuity are verified from Barrick's
official releases. The source corpus supports the operating variables and
valuation interfaces below, but asset ownership, financial perimeter and
structural events remain controlled inputs that must be reconciled before a
final valuation is reported.}

\section{Corporate identity and the choice of case study}

The project retained its historical working name, ``Barrick Gold'', but the
legal issuer is now \emph{Barrick Mining Corporation}. Shareholders approved
the change from Barrick Gold Corporation on 6 May 2025. Trading under the new
name began on 9 May 2025; the New York Stock Exchange ticker changed from
\texttt{GOLD} to \texttt{B}, while the Toronto Stock Exchange ticker remained
\texttt{ABX} \cite{barrick-name-2025}. The distinction matters whenever a
market-price series is stitched across the corporate action.

Barrick describes itself as a global gold and copper exploration, development
and mining company \cite{barrick-about-2026}. Gold remains central, while the
copper portfolio means that Barrick is not a one-factor claim on the gold
price. It is nevertheless a useful case study for this project because gold,
production and cost uncertainty can be connected to operating margin, free
cash flow and a DCF. That connection brings the team's corporate-valuation
interests into the same framework as its stochastic modelling work.

\section{The company perimeter is a model input}

% Sources: GAP-COMPANY-01--04, GAP-DATA-01/03/04; Team 4 MAP-T4-01--12.
A mining-company valuation cannot define the corporate perimeter implicitly.
The perimeter determines which production volumes, costs, taxes, royalties,
cash flows, debt claims and shares belong in the numerator and denominator of
the valuation. It must therefore be frozen at the valuation date and identify,
for every admitted operating asset, the ownership basis, reporting basis,
status and currency. The current Team material supplies a component-driven
method, not an approved current portfolio table.

\gapbox{Populate the dated Barrick corporate and operating perimeter from
approved primary corporate evidence. Legal name and listing continuity are
closed; ownership shares, operating status, structural events and the choice
between consolidated and attributable measures remain open.}

\section{Operating quantities supplied by the source programme}

% Sources: Team 4 MAP-T4-01--12; CONTENT-LEDGER BG-U-10; GAP-VAL-01/02/07.
Team 4 establishes that the operating layer is component-driven rather than a
single revenue-growth extrapolation. Production and cost forecasts are
estimated separately, while a stochastic commodity-price engine supplies the
price input used in the operating distribution. This is sufficient to define
the required contract even though it is insufficient to populate the contract
with current Barrick figures.

\begin{table}[H]
\centering
\caption{Corporate operating contract implied by the Team 4 source module.}
\label{tab:company-operating-contract}
\small
\renewcommand{\arraystretch}{1.12}
\begin{tabularx}{\textwidth}{@{}p{2.70cm}X X@{}}
\toprule
\textbf{Dimension} & \textbf{Required dated input} &
\textbf{Reason it cannot remain implicit} \\
\midrule
Production & Ounces produced and sold by asset, ownership basis, grade,
recovery and forecast frequency. & Volumes determine revenue exposure and
must reconcile consolidated versus attributable reporting. \\
Costs & CoS, cash cost, AISC components, royalties, sustaining and growth
capital with explicit units. & Mixing cash and accrual measures can double
count D\&A, sustaining investment or royalties. \\
Commodity price & Physical-gold and realised-price paths in currency per troy
ounce, including basis, hedging and royalty treatment. & GLD is quoted in USD
per share and cannot enter revenue with an implicit one-to-one conversion. \\
Operating output & Revenue, component margin and the precise EBITDA/EBIT
definition by year and scenario. & The DCF requires a reconciled accounting
quantity, not a label carried over from an operating model. \\
Dependence & Joint or explicitly deterministic relations among price,
production and costs. & Independent broadcasts and untested correlations can
materially change the tails of the cash-flow distribution. \\
\bottomrule
\end{tabularx}
\end{table}

\section{Financial and structural perimeter}

% Sources: Team 5 T5-02/T5-04/T5-10--13; GAP-DATA-03--05;
% GAP-VAL-03/04/06.
The financial perimeter is distinct from the operating perimeter. It fixes the
reporting statements and scale used to derive taxes, reinvestment and capital
cost; it also identifies net debt, lease or debt-equivalent claims, minority
interests, non-operating assets and diluted shares for the equity bridge.
Listing continuity and foreign-exchange treatment matter because a market
comparison is meaningful only when model value and quoted price refer to the
same security, currency and share basis.

\gapbox{Freeze the reporting date, currency and scale; reconcile debt, cash,
minorities, non-operating assets, other claims and diluted shares; document the
security and corporate-action policy used for the contemporaneous market-price
comparison.}

\section{Risk boundary and permitted use}

% Sources: GAP-COMPANY-04; GAP-VAL-02/06; CL-009/013.
Geography, reserve life, execution risk and structural events affect the
operating horizon, the dependence of production and costs, and the credibility
of any terminal assumption. Those facts require current corporate evidence and
cannot be inferred from a market-price series. Until that evidence is admitted,
this chapter defines a valuation contract and its failure conditions; it is not
a substitute company profile and does not support a final target value.
